\section{Discussion}\label{sec:discussion}

\subsection{Forward model or inverse?}

\begin{table}[!htbp]
\centering
\caption{Origin of each difficulty.}
\label{tab:verdict}
\small
\begin{tabular}{p{0.38\linewidth}p{0.22\linewidth}p{0.32\linewidth}}
\toprule
Difficulty & Origin & Remedy\\
\midrule
One scan cannot separate $\Mzero,\Bone,\Tone$ (\cref{cor:onescan}) & forward model (exact) & two scans with a common $\Bone$\\
Small flip angles: leading-order invariance \eqref{eq:sym}, severe ill-conditioning & forward model $+$ protocol & a pulse whose effective angle does not scale with $\Bone$ (near $360^\circ$, as published)\\
Residual $\Tone$ bound ($40.5\,\sigma/\Mzero$): confounding with $\Ttwo$, $\Ttwos$, $\Mzero$, $\Bone$ & forward model $+$ protocol & SNR, averaging, spatial priors; protocol tuning ($\le20\%$ for flip angles)\\
Exact magnitude alias $\alpha_2=360^\circ\mp x$ at physiological $\Bone$ & forward model: magnitudes lose $\sgn\sin\alpha$ & FID-sign test, or complex data with a shared phase\\
Exact complex alias $\alpha_2>540^\circ$ & forward model: $2\pi$ periodicity & range bound $\Bone<1.64$\\
Error 3--7$\times$ the bound; bias and undefined $\Tone$ at moderate SNR & analytic inverse (sequential, closed form) & joint ML or magnitude least squares with the FID sign\\
$\Tone$ bound 57.1 instead of 28.0 with $\Bone$ known; $\Tone$ bias $-2\varepsilon$ per $\Bone$ error $\varepsilon$ & design: scan-1-only $\Tone$ step, low-resolution scan 2 & joint estimation; full-resolution scan 2 if affordable\\
Estimates on $\Rtwop=0$ at low SNR & model domain $+$ small true $\Rtwop$ & report; Bayesian or constrained bounds\\
One start insufficient & nonconvex likelihood & multiple starts (affordable with exact Jacobians)\\
In-vivo calibration $\beta=1.24$ & model mismatch (not variance) & richer forward model (\cref{sec:future})\\
\bottomrule
\end{tabular}
\end{table}

The published protocol is sound from an information standpoint: the near-$360^\circ$ pulse lifts
the leading-order invariance that makes $\Bone$/$\Tone$ separation hopeless with small angles, and the
data contain the information that the bounds promise. The published inverse does not extract
it efficiently; joint estimation does, for complex data and, with one bit of phase, for
magnitudes. The phantom (\cref{sec:phantom}) confirms these conclusions spatially and
quantifies the cost of the low-resolution second scan.

No unbiased voxel-wise estimator can beat the bounds of \cref{sec:information}. Further gains
require prior information: $\Bone$ is spatially smooth and can be treated as nearly known (removing
its factor of about $1.45$), while the confounding with $\Mzero$ and the transverse relaxation times
can only be reduced by priors on the parameter maps themselves. Such estimators are biased by
design, and their bias must be measured rather than ignored.

\subsection{Limitations}

The forward model assumes instantaneous RF pulses, a Lorentzian intravoxel frequency
distribution, no diffusion, no magnetisation transfer, no flow or motion, a single
compartment per voxel, and perfect steady state; the time from the pulse to the first echo
is assumed (4.5\,ms). The CRLB is local and assumes an exact model, Gaussian noise, an interior
true parameter and locally unbiased estimators. The digital phantom is noise-only and
model-consistent: it validates estimators and uncertainty maps but not the model itself.

\subsection{Future work}\label{sec:future}

\paragraph{Model mismatch.} The in-vivo calibration $\beta=1.24$ of~\cite{cheng2019}, unexplained by
the authors and absent in phantoms, indicates that the forward model misses something in
vivo. Candidates, each of which can be added to the simulator and assessed by fitting
mismatched data with the current model: (i) finite hard pulses with off-resonance during the
pulse (their Eq.~8 nutation function) and an inaccurate ratio of the actual flip angles; (ii)
non-Lorentzian intravoxel frequency distributions (e.g.\ Gaussian or susceptibility-driven),
which change $\exp(-\Rtwop|\tau|)$; (iii) diffusion attenuation of the higher-order pathways
under the unbalanced gradient (already available in the EPG simulator); (iv) magnetisation
transfer, which lowers the effective longitudinal magnetisation under the large pulse; (v)
partial volume (two compartments per voxel); (vi) $B_0$ drift or motion between the sequential
scans, which corrupts the relative phase used by the branch test and by complex ML. The
measure of interest is the bias each effect induces in $\Tone$ and $\Bone$ relative to the noise
spread, and whether a single global factor like $\beta$ can absorb it.

\paragraph{Rician maximum likelihood.} Magnitude least squares ignores the Rician noise floor.
At the published protocol the induced bias is small (\cref{sec:magnitude}), but it grows with
lower SNR, smaller flip angles or protocols with weaker pathways. Two routes are natural.
(i) Direct ML with the Rician likelihood
$\log p(m\mid A)=\log m-\log\sigma^2-(m^2+A^2)/(2\sigma^2)+\log I_0(mA/\sigma^2)$, whose gradient and
Gauss--Newton approximation follow from the exact Jacobian of $A=|s(\theta)|$. (ii) An EM algorithm
that treats the lost phase as missing data: with $\kappa=m|s_t|/\sigma^2$, the E-step gives
$\mathbb E[|z|\mid m]=m\,I_1(\kappa)/I_0(\kappa)$ and the M-step is magnitude least squares on these
shrunk pseudo-data, which is the existing fit. Either should be combined with the FID-sign
constraint, which the alias of \cref{prop:equiv} requires for any magnitude likelihood.
The Rician Fisher information of \cref{sec:rician} provides the reference bound.

\paragraph{Hierarchical (empirical-Bayes) estimation.} A hierarchical model with tissue classes
(a Gaussian mixture in $(\log\Tone,\log\Ttwo,\log\Ttwos)$) and a smooth parametric $\Bone$ field, whose
hyperparameters are estimated by EM with a Laplace E-step (MAP fit plus
$(J^{\mathsf T}J+\Sigma_k^{-1})^{-1}$ per voxel and class), would use the subject's own data as the
prior. It targets exactly the weak direction of \cref{fig:eigen} and makes the $\Bone$ field
effectively known; its bias would be evaluated on the phantom, including the lesion.

\paragraph{Learned estimators.} Networks trained on simulated signals~\cite{cheng2020} are
biased estimators whose prior is the training distribution. The tools here---exact bounds,
efficient ML as a reference, calibrated uncertainty, and the phantom---are what is needed to
measure their bias and variance against the CRLB rather than only their average error.

\paragraph{Acquisition.} Direct CRLB-based optimisation of $\TR$, flip angles, echo placement
and scan-2 coverage, now cheap (\cref{sec:jacobian}); and the joint use of a low-resolution
scan 2 inside a single likelihood (modelling its $k$-space truncation) instead of the two-stage
pipeline.
